\documentclass{amsart}

% 数学相关宏包
\usepackage{amsmath,mathtools,amsthm,amssymb}
\usepackage{bbm}
\usepackage{mathrsfs}
\usepackage{centernot}

% 图形和绘图
\usepackage{graphicx,float}
\usepackage{tikz,tikz-cd}
\usepackage{pgfplots}
\usepackage[framemethod=TikZ]{mdframed}
\usetikzlibrary{decorations.pathmorphing,positioning,arrows,automata,calc,shapes,patterns}
\pgfplotsset{compat=1.18}
\usepgfplotslibrary{statistics}

% 表格和列表
\usepackage{booktabs,multirow,colortbl}
\usepackage{enumitem}
\setlist[itemize,enumerate,description]{noitemsep}
\setlistdepth{9}

% 算法
\usepackage{algorithm,algorithmic}

% 字体和颜色
\usepackage{xcolor,listings}
\usepackage{xeCJK}
\setCJKmainfont{Noto Serif CJK SC}

% 页面和链接
\usepackage{fancyhdr,caption}
\usepackage{hyperref}
\usepackage{cleveref}

% 工具宏包
\usepackage{etoolbox,letltxmacro,docmute}
\usepackage{gensymb}

% 自定义设置
\renewcommand{\arraystretch}{1.5}
\let\originalmathbb\mathbb
\renewcommand{\mathbb}[1]{%
    \ifstrequal{#1}{1}{\mathbbm{#1}}{\originalmathbb{#1}}%
}

% 颜色定义
\definecolor{lightblue}{RGB}{173,216,230}
\definecolor{lightgreen}{RGB}{144,238,144}
\definecolor{lightyellow}{RGB}{255,255,224}
\definecolor{lightgray}{RGB}{211,211,211}

		% 超链接设置
		\hypersetup{
		    colorlinks=true,
		    linkcolor=red,
		    filecolor=magenta,
		    urlcolor=cyan,
		    pdftitle={\@title},
		    pdfauthor={\@author},
		    pdfsubject={数学笔记},
		    pdfkeywords={数学},
		    bookmarksnumbered=true,
		    bookmarksopen=true,
		    bookmarksopenlevel=6,
		    pdfstartview=FitH,
		    pdfpagemode=UseOutlines,
		    pdfborder={0 0 0}
		}

% 定理环境
\theoremstyle{plain}
\newtheorem{theorem}{Theorem}[section]
\newtheorem{lemma}[theorem]{Lemma}
\newtheorem{corollary}[theorem]{Corollary}
\newtheorem{proposition}[theorem]{Proposition}

\theoremstyle{definition}
\newtheorem{definition}[theorem]{Definition}
\newtheorem{example}[theorem]{Example}
\newtheorem{exercise}[theorem]{Exercise}
\newtheorem{problem}[theorem]{Problem}

\theoremstyle{remark}
\newtheorem{remark}[theorem]{Remark}
\newtheorem{note}[theorem]{Note}
\newtheorem{observation}[theorem]{Observation}

% 解答环境
\newtheorem*{solution}{Solution}
\newtheorem*{proof*}{Proof}

% 图片路径
\graphicspath{{F:/资料/2025-summer/figures/}}

\author{乐绎华, Yue Yihua}
\address{中山大学, Sun Yat-sen University}
\email{yueyh@mail2.sysu.edu.cn}
\date{\today}
\title{Mathematical Notes}

\begin{document}

\maketitle
\tableofcontents

The symbol \lstinline{⊢} means "we need to prove the following".

\begin{table}[h]
	\centering
	\begin{tabular}{|c|c|c|c|}
		\hline
		策略 & 主要作用 & 何时使用 & 行为特点 \\
		\hline
		\lstinline{rw} & \textbf{重写}表达式的局部或整体 & 当你想用一个等式替换目标中的一部分时 & 替换匹配的子表达式，不改变目标的主要结构 \\
		\hline
		\lstinline{exact} & \textbf{直接完成}目标 & 当你已经有了一个与当前目标类型\textbf{完全匹配}的证明项时 & 不产生新的子目标，直接关闭当前目标 \\
		\hline
		\lstinline{apply} & \textbf{应用}一个函数、定理或引理来分解目标 & 当你想用一个形如 \lstinline{A → B} 的表达式来证明 \lstinline{B} 并生成 \lstinline{A} 作为子目标时 & 尝试匹配结论，然后将前提作为新的子目标（可能多个） \\
		\hline
	\end{tabular}
\end{table}
\lstinline{norm_num} 将涉及常量的数值表达式简化为其最简形式.

Task: Lean 的证明和命题结构围绕着\textbf{类型论}和**战术（tactics）**展开。理解这些基本结构是编写 Lean 证明的关键。

\subsection{Lean 证明的基本结构}

Lean 中的证明通常是\textbf{构造性的}。这意味着，要证明一个命题，你需要构造一个“证据”或“项”来满足这个命题的类型。

\begin{enumerate}
	\item \textbf{\lstinline{example} 或 \lstinline{theorem}：引入一个要证明的命题}
	\begin{itemize}
		\item \textbf{\lstinline{example}}: 用于引入一个\textbf{不带名称}的命题或定义。常用于教学示例或你不需要在其他地方引用的临时证明。
	\end{itemize}
\begin{lstlisting}[language=lean]
example : 2 + 2 = 4 := by rfl
\end{lstlisting}	\begin{itemize}
		\item \textbf{\lstinline{theorem}}: 用于引入一个\textbf{带名称}的命题或定理。这是你在数学库中定义可重用结果的标准方式。
	\end{itemize}
\begin{lstlisting}[language=lean]
theorem two_plus_two_eq_four : 2 + 2 = 4 := by rfl
\end{lstlisting}\lstinline{theorem} 后的名称 (\lstinline{two_plus_two_eq_four}) 允许你在其他证明中引用它。
	\item \textbf{\lstinline{:}：指定命题的类型}
	\begin{itemize}
		\item 在 \lstinline{example} 或 \lstinline{theorem} 之后，冒号 \lstinline{:} 用于指定你要证明的命题本身。在 Lean 中，命题的类型通常是 \lstinline{Prop}。
	\end{itemize}
\begin{lstlisting}[language=lean]
example : 2 + 2 = 4  -- 2 + 2 = 4 是一个 Prop 类型的命题
\end{lstlisting}	\item \textbf{\lstinline{:= by}：开始战术证明块}
	\begin{itemize}
		\item \lstinline{:=} 表示定义或等同于。
		\item \lstinline{by} 引入一个\textbf{战术块}。这意味着你将使用一系列 Lean 的战术来逐步构建证明，而不是直接写一个完整的证明项。这是 Lean 中最常见的证明风格。
	\end{itemize}
\begin{lstlisting}[language=lean]
example : 2 + 2 = 4 := by
  -- 在这里写战术
\end{lstlisting}	\item \textbf{战术（Tactics）：构建证明的步骤}
战术是 Lean 用来改变或解决证明目标（goal）的命令。它们是证明的核心。
	\begin{itemize}
		\item \textbf{\lstinline{rfl} (reflexivity)}: 证明一个表达式等于自身。常用于简单的等式，如 \lstinline{2 + 2 = 4}，当两边经过计算后是相同的。
	\end{itemize}
\begin{lstlisting}[language=lean]
example : 2 + 2 = 4 := by rfl
\end{lstlisting}	\begin{itemize}
		\item \textbf{\lstinline{exact <term>}}: 当你有一个与当前目标\textbf{完全匹配}的证明项时，直接提供它来完成目标。
	\end{itemize}
\begin{lstlisting}[language=lean]
example (h : P) : P := by exact h
\end{lstlisting}	\begin{itemize}
		\item \textbf{\lstinline{apply <lemma/theorem/function>}}: 应用一个函数、引理或定理。它会尝试将当前目标与应用项的结论部分匹配，然后将应用项的前提部分变成新的子目标。
	\end{itemize}
\begin{lstlisting}[language=lean]
example (P Q : Prop) (hpq : P → Q) (hp : P) : Q := by
  apply hpq -- 目标从 Q 变为 P
  exact hp  -- 目标 P 被 hp 解决
\end{lstlisting}	\begin{itemize}
		\item \textbf{\lstinline{intro <name>} / \lstinline{intros <names>}}: 如果目标是形如 \lstinline{A → B} 的蕴含，\lstinline{intro} 会引入一个假设，将目标从 \lstinline{A → B} 变为 \lstinline{B}，并将 \lstinline{A} 命名为一个局部假设。
	\end{itemize}
\begin{lstlisting}[language=lean]
example (P Q : Prop) : P → Q → P := by
  intro hP -- 引入 P 的假设，命名为 hP。目标变为 Q → P
  intro hQ -- 引入 Q 的假设，命名为 hQ。目标变为 P
  exact hP -- 解决目标
\end{lstlisting}	\begin{itemize}
		\item \textbf{\lstinline{have <name> : <type> := by <proof>} / \lstinline{have <name> : <type> := <term>}}: 引入一个\textbf{中间结论或引理}。你可以先证明这个中间结论，然后在其后的证明中像使用假设一样使用它。
	\end{itemize}
\begin{lstlisting}[language=lean]
example (a b c : Nat) (h₁ : a = b) (h₂ : b = c) : a = c := by
  have h₃ : a = c := Eq.trans h₁ h₂ -- 使用 Eq.trans 证明 a = c
  exact h₃
\end{lstlisting}或者使用战术证明 \lstinline{have} 的部分：
\begin{lstlisting}[language=lean]
example (a b c : Nat) (h₁ : a = b) (h₂ : b = c) : a = c := by
  have h₃ : a = c := by rw [h₁, h₂]
  exact h₃
\end{lstlisting}	\begin{itemize}
		\item \textbf{\lstinline{rw <equation>} (rewrite)}: 根据一个等式或双向蕴含来重写目标或假设中的一部分。
	\end{itemize}
\begin{lstlisting}[language=lean]
example (a b : Nat) (h : a = b + 1) : a + 2 = b + 1 + 2 := by
  rw h -- 将目标中的 'a' 替换为 'b + 1'
  rfl  -- 目标变为 b + 1 + 2 = b + 1 + 2，用 rfl 解决
\end{lstlisting}	\begin{itemize}
		\item \textbf{\lstinline{use <term>, <proofs>}}: 用于证明\textbf{存在量词} \lstinline{∃ x, P x}。你需要提供一个具体的项 \lstinline{x} 以及证明 \lstinline{P x} 成立的证据。
	\end{itemize}
\begin{lstlisting}[language=lean]
example : ∃ x : Nat, x > 10 := by
  use 11 -- 提供 x = 11
  norm_num -- 证明 11 > 10
\end{lstlisting}	\begin{itemize}
		\item \textbf{\lstinline{norm_num}}: 规范化数值计算（如前所述）。它能自动处理常量表达式的算术和比较。
	\end{itemize}
\begin{lstlisting}[language=lean]
example : 2 * (3 + 4) = 14 := by norm_num
\end{lstlisting}\end{enumerate}

\subsection{你提供的例子：\lstinline{example : ∃ x : ℝ, 2 < x ∧ x < 3 := by ...}}

这是一个证明\textbf{存在量词}的经典例子。我们来细分它：

\begin{lstlisting}[language=lean]
example : ∃ x : ℝ, 2 < x ∧ x < 3 := by
  -- 目标是：存在一个实数 x，使得 2 < x 且 x < 3。
  -- 这是一个存在性证明。你需要找到一个具体的 x，并证明它满足条件。

  have h1 : 2 < (5 : ℝ) / 2 := by
    norm_num -- 使用 norm_num 证明 2 < 2.5 成立。
    -- (5 : ℝ) / 2 明确指出 5 和 2 应该被视为实数进行除法，结果是 2.5。

  have h2 : (5 : ℝ) / 2 < 3 := by
    norm_num -- 使用 norm_num 证明 2.5 < 3 成立。

  use 5 / 2, h1, h2
  -- `use` 战术用于解决存在量词。
  -- 它需要三部分：
  -- 1. `5 / 2`: 具体的存在项 x (在这里是 2.5)。
  -- 2. `h1`: 证明 2 < x 的证据（即 2 < 2.5）。
  -- 3. `h2`: 证明 x < 3 的证据（即 2.5 < 3）。
  -- 由于 2 < x 和 x < 3 是一个合取 (`∧`)，`use` 需要分别提供这两个证据。
\end{lstlisting}
这个例子展示了 Lean 证明的几个核心概念：

\begin{itemize}
	\item \textbf{存在性证明}需要你提供一个具体的例子。
	\item \textbf{中间步骤}可以使用 \lstinline{have} 引入并证明。
	\item \textbf{自动化战术}如 \lstinline{norm_num} 可以极大地简化数值相关的证明。
	\item \textbf{类型注解} (\lstinline{(5 : ℝ) / 2}) 在处理不同数值类型时很重要，确保 Lean 理解你在哪个域（例如，整数 \lstinline{ℕ}，整数 \lstinline{ℤ}，有理数 \lstinline{ℚ}，实数 \lstinline{ℝ}）中进行计算。
\end{itemize}

掌握这些基本结构和常用战术是开始在 Lean 中进行形式化证明的良好起点。你还有哪些想了解的特定证明类型或概念吗？


\end{document}